\documentclass[11pt]{article}
\usepackage[T1]{fontenc}
\usepackage{lmodern}
\usepackage[margin=1in]{geometry}
\usepackage{amsmath,amssymb,amsthm,mathtools}
\usepackage{microtype}
\usepackage{needspace}
\usepackage[dvipsnames]{xcolor}
\usepackage{hyperref}
\hypersetup{colorlinks=true,linkcolor=MidnightBlue,citecolor=MidnightBlue,
  urlcolor=MidnightBlue,pdftitle={Scalar compressions and asymptotically optimal tough error models},
  pdfsubject={Review draft; independent proof and prior-work review pending}}
\newcommand{\C}{\mathbb C}
\newcommand{\E}{\mathcal E}
\newcommand{\V}{\mathcal V}
\DeclareMathOperator{\Tr}{Tr}
\DeclareMathOperator{\rank}{rank}
\DeclareMathOperator{\ran}{ran}
\DeclareMathOperator{\Span}{span}
\newtheorem{theorem}{Theorem}[section]
\newtheorem{proposition}[theorem]{Proposition}
\newtheorem{corollary}[theorem]{Corollary}
\theoremstyle{remark}
\newtheorem{remark}[theorem]{Remark}
\numberwithin{equation}{section}
\setlength{\emergencystretch}{1em}
\title{Scalar compressions and asymptotically optimal\protect\\ tough error models}
\author{}
\date{13 November 2026\\[4pt]
  \small Review draft: independent proof and prior-work review pending}

\begin{document}
\maketitle
\begin{abstract}
For an $e$-dimensional complex space of errors acting on $\C^n$, let
$c(e,n)$ be the largest code dimension guaranteed by the Knill--Laflamme
conditions without further assumptions on the errors. We prove
\[
 \max\!\left\{1,\left\lfloor\frac{n+1}{e^2+1}\right\rfloor\right\}
 \le c(e,n)\le
 \left\lceil\frac{n}{e^2+1}\right\rceil
 \qquad (1\le e\le n^2).
\]
The two bounds differ by at most one. Explicit error spaces attain the
upper bound, so their largest codes are within one dimension of the
worst possible value. The lower bound follows from a scalar-compression
theorem: every $D$-dimensional operator system in $M_n(\C)$ has a rank-$k$
scalar compression when $2\le k\le n$ and $n\ge Dk-1$.
We minimize the squared traceless parts of the compressed observables.
At a positive minimum, stationarity supplies two linear dependencies;
these leave a common orthogonal direction along which a sum of second
variations is negative. Thus $c(e,n)=n/(e^2+1)+O(1)$ uniformly in $e,n$,
with relative asymptotic equality whenever $n/(e^2+1)\to\infty$.
Some finite integer values remain undetermined.
\end{abstract}

\section{The guaranteed code dimension}

How large an exactly correcting subspace must exist if only the dimension
of the error space is known? This is the tough-error-model problem in
Kr\"uger and Werner's collection~\cite[Problem~14]{KW05} and in its
current catalog formulation~\cite{QIQC}.
Let $\E\subseteq M_n(\C)$ be a complex linear space with
$\dim_{\C}\E=e$. A nonzero subspace $C\subseteq\C^n$ corrects $\E$ if
its orthogonal projection $P_C$ satisfies
\begin{equation}\label{eq:kl}
 P_C A^*B P_C=\lambda(A,B)P_C\qquad(A,B\in\E)
\end{equation}
for suitable scalars $\lambda(A,B)$. This is the algebraic
Knill--Laflamme condition~\cite{KL97,KLV00}. Define
\begin{equation}\label{eq:minimax}
 \alpha(\E)=\max\{\dim C:C\text{ satisfies \eqref{eq:kl}}\},
 \qquad
 c(e,n)=\min_{\substack{\E\subseteq M_n(\C)\\\dim_{\C}\E=e}}
 \alpha(\E).
\end{equation}
Every line satisfies \eqref{eq:kl}, so these quantities are positive
integers. The definition imposes neither $I_n\in\E$ nor a
trace-preserving normalization on a basis of $\E$. In particular, $e$
counts a dimension of an arbitrary error space, rather than the Kraus
rank of a normalized channel.

\Needspace{7\baselineskip}
\begin{theorem}\label{thm:main}
For all integers $n\ge1$ and $1\le e\le n^2$,
\begin{equation}\label{eq:main}
 L(e,n):=\max\!\left\{1,\left\lfloor\frac{n+1}{e^2+1}\right\rfloor\right\}
 \le c(e,n)\le
 U(e,n):=\left\lceil\frac{n}{e^2+1}\right\rceil.
\end{equation}
There is an explicit $e$-dimensional error space $\E_{e,n}$ with
$\alpha(\E_{e,n})=U(e,n)$. The gap $U(e,n)-L(e,n)$ is at most one.
\end{theorem}

Thus the construction is asymptotically optimal whenever the guaranteed
code dimension diverges. More precisely,
\begin{equation}\label{eq:asymptotic}
 \left|c(e,n)-\frac{n}{e^2+1}\right|<1,
 \qquad
 \frac{c(e,n)}{n/(e^2+1)}\longrightarrow1
 \quad\text{if }\frac{n}{e^2+1}\longrightarrow\infty.
\end{equation}
The error bound is absolute and allows $e$ to vary with $n$. This answers
the asymptotically matching-bounds version of the problem. It does not
determine the exact rounding at every integer pair.

The new lower bound comes from treating all scalar-compression constraints
at once. For a Hermitian basis $I_n,H_1,\ldots,H_{D-1}$ of an operator
system, consider the sum of the squared traceless parts of
$X^*H_jX$, where $X$ is an $n\times k$ isometry. If this sum is positive
at a minimum, stationarity makes the columns of
$[X,H_1X,\ldots,H_{D-1}X]$ linearly dependent in at least two directions.
When $n\ge Dk-1$, a vector perpendicular to all these columns remains.
Rotating each column of $X$ toward that vector gives a negative sum of
second derivatives, contradicting minimality. Section~\ref{sec:variance}
gives the argument without assumptions on the spectra or commutators of
the observables.

For comparison, the historical estimate recorded in~\cite{KW05}, from
the Knill--Laflamme--Viola construction~\cite{KLV00}, has scale
$n/[e^2(e^2+1)]$. Reimpell also studied this problem through repeated
Hermitian compression and numerical optimization~\cite[Section~3.3.4]{Reimpell}.
In operator-system language, a scalar compression is a quantum anticlique.
Weaver's general existence argument~\cite[Theorem~2.10]{Weaver} uses
diagonal compression followed by a convex partition, with an ambient
dimension requirement quadratic in the operator-system dimension.
Allen and Kornell~\cite[Theorem~2.7]{AK} obtain sharper bounds in small
dimensions: dimension at most $2\ell+1$ guarantees rank $k$ when
$n>3^\ell(k-1)$, and dimension at most $2\ell+2$ when
$n>2\cdot3^\ell(k-1)$. In particular, their result gives
$c(2,n)\ge\lceil n/9\rceil$, whereas \eqref{eq:main} has asymptotic
coefficient $1/5$. These comparisons concern the stated general bounds;
special low-dimensional results can have better integer thresholds than
the theorem below.

The upper construction uses standard rank-one spanning vectors from
Li and Tam~\cite[Section~2.1]{LiTam} and the block-rank obstruction
in Weaver~\cite[Proposition~2.1]{Weaver}. Section~\ref{sec:models}
assembles these ingredients into an $e$-dimensional error space, including
the zero label required by the unnormalized formulation. The construction
and its exact largest-code calculation are included in full. Geometric
and variational descriptions of scalar compression also occur in recent
work on Pauli detection~\cite{Gupta,Sun}; the theorem here is a uniform
existence statement for arbitrary operator systems, with no prescribed
Pauli or representation structure.

\section{A variational scalar-compression theorem}\label{sec:variance}

An \emph{operator system} $\V\subseteq M_n(\C)$ is a complex linear
subspace containing $I_n$ and satisfying $\V^*=\V$. A rank-$k$
\emph{scalar compression} of $\V$ is an isometry
$X:\C^k\longrightarrow\C^n$ for which $X^*HX$ is a scalar multiple
of $I_k$ for every $H\in\V$. Equivalently, $P=XX^*$ is a rank-$k$
projection with $P\V P=\C P$.

\Needspace{10\baselineskip}
\begin{theorem}\label{thm:compression}
Let $\V\subseteq M_n(\C)$ be an operator system of complex dimension
$D$. If $2\le k\le n$ and
\begin{equation}\label{eq:threshold}
 n\ge Dk-1,
\end{equation}
then $\V$ has a rank-$k$ scalar compression. Consequently, every such
system has a scalar compression of rank at least
\begin{equation}\label{eq:rank-bound}
 \min\!\left\{n,\max\!\left\{1,
 \left\lfloor\frac{n+1}{D}\right\rfloor\right\}\right\}.
\end{equation}
For $D=1$, the entire space is a scalar compression.
\end{theorem}

\begin{proof}
The self-adjoint part of $\V$ has real dimension $D$ and spans $\V$
over $\C$. Choose a Hermitian basis $I_n,H_1,\ldots,H_{D-1}$.
For an isometry $X\in M_{n\times k}(\C)$, set
\begin{equation}\label{eq:objective}
 G_j=X^*H_jX,\qquad
 \mu_j=\frac{\Tr G_j}{k},\qquad
 T_j=G_j-\mu_j I_k,\qquad
 F(X)=\sum_{j=1}^{D-1}\Tr(T_j^2).
\end{equation}
Each $T_j$ is Hermitian and traceless. Thus $F(X)\ge0$, with equality
precisely when $X$ is a scalar compression. The function is continuous
on the compact space of isometries and invariant under $X\mapsto XU$
for unitary $U\in M_k(\C)$. Fix a global minimizer $X$ and suppose
that $F(X)>0$.

For an arbitrary variation $Z$ of $X$, differentiation gives
\begin{equation}\label{eq:differential}
 dF_X[Z]=4\operatorname{Re}\Tr\!\left[
 \left(\sum_jH_jXT_j\right)^*Z\right].
\end{equation}
Every $Z$ satisfying $X^*Z=0$ is tangent to the isometries at $X$.
Since \eqref{eq:differential} vanishes for all such $Z$, the horizontal
component of the gradient is zero. Hence
\begin{equation}\label{eq:stationarity}
 \sum_j H_jXT_j=XS,\qquad
 S=\sum_jG_jT_j.
\end{equation}
Each $G_j$ commutes with $T_j$, so $S$ is Hermitian. Introduce the block
matrices
\begin{equation}\label{eq:dependencies}
 M=\begin{bmatrix}X&H_1X&\cdots&H_{D-1}X\end{bmatrix},
 \qquad
 B=\begin{bmatrix}-S\\T_1\\\vdots\\T_{D-1}\end{bmatrix}.
\end{equation}
Equation~\eqref{eq:stationarity} says $MB=0$. Since $F(X)>0$, at least
one $T_j$ is a nonzero traceless Hermitian matrix. Such a matrix has
both a positive and a negative eigenvalue, and therefore rank at least
two. As a block of $B$, it implies $\rank B\ge2$. The matrix $M$ has
$Dk$ columns, and its kernel contains the range of $B$. Consequently,
\begin{equation}\label{eq:rank-defect}
 \rank M\le Dk-2<n.
\end{equation}
Choose a unit vector $y\in\C^n$ orthogonal to the range of $M$. Then
\begin{equation}\label{eq:orthogonality}
 X^*y=0,\qquad X^*H_jy=0\quad(1\le j\le D-1),
\end{equation}
where the second equality uses $H_j^*=H_j$.

Write $X=[x_1,\ldots,x_k]$. For each $a\in\{1,\ldots,k\}$, replace
the column $x_a$ by $\cos\theta\,x_a+\sin\theta\,y$, leaving all
other columns fixed. Denote this isometry by $X_a(\theta)$ and put
$F_a(\theta)=F(X_a(\theta))$. Let $f_a$ be the $a$th coordinate vector
of $\C^k$, $E_a=f_af_a^*$, and
$C_a(\theta)=I_k+(\cos\theta-1)E_a$. Then
\[
 X_a(\theta)=XC_a(\theta)+\sin\theta\,yf_a^*.
\]
Writing $h_j=y^*H_jy$, the orthogonality in
\eqref{eq:orthogonality} gives the exact identity
\begin{equation}\label{eq:rotation-exact}
 X_a(\theta)^*H_jX_a(\theta)
 =C_a(\theta)G_jC_a(\theta)+\sin^2\theta\,h_jE_a.
\end{equation}
Its expansion at zero is
\begin{equation}\label{eq:rotation-expansion}
 G_j+\theta^2\left(h_jE_a-\tfrac12(E_aG_j+G_jE_a)\right)
 +O(\theta^4).
\end{equation}
There is no first-order term. Since
$\Tr(T_j^2)=\Tr(G_j^2)-(\Tr G_j)^2/k$, differentiation yields
\begin{align}
 F_a''(0)
 &=4\sum_j\Tr\!\left[T_j
 \left(h_jE_a-\tfrac12(E_aG_j+G_jE_a)\right)\right]\notag\\
 &=4\sum_j\left(h_j(T_j)_{aa}-(G_jT_j)_{aa}\right).
 \label{eq:second-variation}
\end{align}
The last step again uses $G_jT_j=T_jG_j$. Summing over the $k$ column
rotations cancels the terms containing $h_j$, because $\Tr T_j=0$.
Moreover, $\Tr(G_jT_j)=\Tr(T_j^2)$. Therefore
\begin{equation}\label{eq:negative-sum}
 \sum_{a=1}^kF_a''(0)
 =-4\sum_j\Tr(T_j^2)=-4F(X)<0.
\end{equation}
Every $F_a''(0)$ must be nonnegative at a minimum, which is a
contradiction. Thus $F(X)=0$.

A rank-one scalar compression always exists. If $D\ge2$ and the rank
in \eqref{eq:rank-bound} is at least two, it satisfies
\eqref{eq:threshold} and is at most $n$, so the first part applies.
If $D=1$, then $\V=\C I_n$ and rank $n$ is immediate.
\end{proof}

\begin{remark}
The same argument shows that any positive stationary value of $F$ has
a direction of negative second variation under \eqref{eq:threshold}.
The existence proof uses compactness to select a global minimum; it
does not provide a running-time guarantee for finding a correcting code.
\end{remark}

\begin{corollary}\label{cor:lower}
Every $e$-dimensional complex error space in $M_n(\C)$ has a correcting
code of dimension at least $L(e,n)$ in \eqref{eq:main}.
\end{corollary}

\begin{proof}
Choose a complex basis $A_1,\ldots,A_e$ of $\E$ and form
\begin{equation}\label{eq:error-system}
 \V_{\E}=\Span_{\C}\{I_n,A_i^*A_j:1\le i,j\le e\}.
\end{equation}
This is an operator system of dimension $d\le e^2+1$.
Let $k=\lfloor(n+1)/(e^2+1)\rfloor$. If $k\ge2$, then
$k\le n$ and $n\ge(e^2+1)k-1\ge dk-1$, so
Theorem~\ref{thm:compression} gives a rank-$k$ scalar compression of
$\V_{\E}$. The pairwise products $A_i^*A_j$ then satisfy
\eqref{eq:kl}, and sesquilinearity extends the condition to all
$A,B\in\E$. When $k\le1$, any line is a correcting code.
\end{proof}

\section{Explicit tough error spaces}\label{sec:models}

We next construct the error spaces attaining the upper bound in
Theorem~\ref{thm:main}. The basic obstruction is the following
block-rank calculation~\cite[Proposition~2.1]{Weaver}.

\Needspace{7\baselineskip}
\begin{proposition}\label{prop:blocks}
Let $P_1,\ldots,P_t$ be nonzero mutually orthogonal projections summing
to $I_n$, and let $\V=\Span_{\C}\{P_1,\ldots,P_t\}$. The largest
rank of a scalar compression of $\V$ is
$\max_s\rank P_s$.
\end{proposition}

\begin{proof}
If a rank-$k$ projection $P$ satisfies $PP_sP=\lambda_sP$ for all $s$,
positivity gives $\lambda_s\ge0$ and summing gives $\sum_s\lambda_s=1$.
For some $s$, $\lambda_s>0$, and hence
$k=\rank(PP_sP)\le\rank P_s$. Conversely, the range of any $P_s$
is a scalar compression of $\V$.
\end{proof}

Assume first that $1\le e<n$, and put $t=\min\{n,e^2+1\}$.
Let $u_1,\ldots,u_e$ be the coordinate vectors in $\C^e$. Make an
ordered list of $e^2$ vectors by taking all $u_i$ first, followed, in
lexicographic order on $i<j$, by
\begin{equation}\label{eq:labels}
 u_i+u_j,\qquad u_i+\mathrm{i}u_j.
\end{equation}
These are the coordinate and pair vectors in the rank-one spanning
construction of Li and Tam~\cite[Section~2.1]{LiTam}, without normalizing
their lengths. Let $v_0=0$, and let $v_1,\ldots,v_{t-1}$ be the first
$t-1$ vectors of the list. Because $t-1\ge e$, all coordinate vectors
are included.

Decompose $\C^n$ into $t$ nonzero coordinate blocks with dimensions
differing by at most one, and denote their projections by
$P_0,\ldots,P_{t-1}$. Explicitly, writing $n=qt+r$ with $0\le r<t$,
take $r$ blocks of dimension $q+1$ and $t-r$ blocks of dimension $q$.
Define diagonal errors
\begin{equation}\label{eq:errors}
 A_a=\sum_{s=0}^{t-1}(v_s)_aP_s\quad(1\le a\le e),
 \qquad \E_{e,n}=\Span_{\C}\{A_1,\ldots,A_e\}.
\end{equation}
On the block labelled by $u_i$, $A_a$ equals $\delta_{ai}I$.
Thus the $A_a$ are linearly independent and $\dim\E_{e,n}=e$.

To identify their product space, write $D_i$ for the projection of the
block labelled by $u_i$. For $i<j$, let $R_{ij}$ and $S_{ij}$ be the
projections of the blocks labelled by $u_i+u_j$ and
$u_i+\mathrm{i}u_j$, respectively; a missing block means the zero
projection. The labels give
\begin{equation}\label{eq:cross-products}
 A_i^*A_j=R_{ij}+\mathrm{i}S_{ij},\qquad
 R_{ij}=\frac{A_i^*A_j+A_j^*A_i}{2},\qquad
 S_{ij}=\frac{A_i^*A_j-A_j^*A_i}{2\mathrm{i}}.
\end{equation}
They also give
\begin{align}
 D_i&=A_i^*A_i-
 \sum_{\substack{p<q\\i\in\{p,q\}}}(R_{pq}+S_{pq}),
 \label{eq:coordinate-blocks}\\
 P_0&=I_n-\sum_iD_i-\sum_{i<j}(R_{ij}+S_{ij}).
 \label{eq:zero-block}
\end{align}
Every block projection therefore belongs to $\V_{\E_{e,n}}$.
Conversely, the errors and all their products are constant on each
block. We have proved
\begin{equation}\label{eq:block-system}
 \V_{\E_{e,n}}=\Span_{\C}\{P_0,\ldots,P_{t-1}\}.
\end{equation}
Proposition~\ref{prop:blocks} now gives
\begin{equation}\label{eq:largest-code}
 \alpha(\E_{e,n})=\left\lceil\frac nt\right\rceil
 =\left\lceil\frac{n}{e^2+1}\right\rceil.
\end{equation}
The last equality holds because $t=e^2+1$ when $n\ge e^2+1$, and
both ceilings equal one otherwise. The rank argument applies to every
subspace, including codes with support across several blocks.

For $e\ge n$, let $\E_{e,n}$ contain all $n$ diagonal matrix units
and extend them by off-diagonal matrix units, in lexicographic order,
to a linearly independent set of $e$ errors. Since $\V_{\E_{e,n}}$
contains the $n$ rank-one coordinate projections, the same rank argument
forces every code to have dimension one. Here
$\lceil n/(e^2+1)\rceil=1$. This covers the remaining parameters,
including $n=e=1$, and proves the upper bound in
Theorem~\ref{thm:main}.

The smallest two-error example displaying all five labels is
\begin{equation}\label{eq:five-block-example}
 A_1=\operatorname{diag}(0,1,0,1,1),\qquad
 A_2=\operatorname{diag}(0,0,1,1,\mathrm{i}).
\end{equation}
Its product system, with the identity adjoined as in
\eqref{eq:error-system}, is the full diagonal algebra in $M_5(\C)$,
so its largest code has dimension one. Replacing the five coordinates
by balanced blocks gives largest code dimension $\lceil n/5\rceil$
for every $n\ge5$.

\section{Asymptotics and the remaining integer gap}

It remains to check the one-dimension gap in Theorem~\ref{thm:main}.
Write $D=e^2+1$ and $n=qD+r$, with $0\le r<D$.
If $r=0$ or $r=D-1$, the two sides of \eqref{eq:main} agree.
For $1\le r\le D-2$ and $q\ge1$, they are $q$ and $q+1$.
If $q=0$, both bounds are one. This proves the stated gap and
\eqref{eq:asymptotic}.
It also gives, uniformly in all admissible parameters,
\begin{equation}\label{eq:uniform-scale}
 c(e,n)=\Theta\!\left(\max\!\left\{1,\frac{n}{e^2+1}\right\}\right).
\end{equation}
For example, absolute comparison constants $1/2$ and $2$ follow
directly from \eqref{eq:main}.

The bounds recover $c(1,n)=\lceil n/2\rceil$, also a consequence of
the Hermitian higher-rank numerical-range formula~\cite{CKZ}.
They give $c(e,n)=1$ whenever $n\le e^2+1$, and exact values at
$n=q(e^2+1)$ and $n=q(e^2+1)-1$ for every admissible positive $q$.
For two errors, this yields $c(2,9)=c(2,10)=2$ and
$c(2,14)=c(2,15)=3$. At $(e,n)=(2,6)$, the present argument leaves
$1\le c(2,6)\le2$.

The corresponding integer question for arbitrary operator systems is
whether every $D$-dimensional system has a scalar compression of rank
$\lceil n/D\rceil$. This is a form of Weaver's conjecture on the lower
quantum Tur\'an numbers~\cite[following Theorem~2.11]{Weaver}.
For $D=3$, the stronger threshold $n\ge3k-2$ is known~\cite[Lemma~2.3]{AK}.
Theorem~\ref{thm:compression} gives $n\ge Dk-1$ for every $D$ but does
not reach the conjectured threshold $n\ge D(k-1)+1$ in general.
The remaining rounding problem for error spaces may also benefit from
their factorization as products $A_i^*A_j$, which was not used beyond
the dimension bound in Corollary~\ref{cor:lower}.

Finally, normalization changes the extremal problem. If a chosen error
basis satisfies $\sum_i A_i^*A_i=I_n$, then the identity already lies
in the product span and $\dim\V_{\E}\le e^2$. The scalar-compression
theorem still applies with that smaller dimension. The explicit models
in Section~\ref{sec:models}, however, use a common zero block and need
not satisfy this normalization. Equation~\eqref{eq:main} concerns the
arbitrary linear error spaces in \eqref{eq:minimax}.

\begin{thebibliography}{99}
\raggedright
\bibitem{KW05}
O.~Kr\"uger and R.~F.~Werner (eds.),
\emph{Some Open Problems in Quantum Information Theory} (2005),
Problem~14, \href{https://arxiv.org/abs/quant-ph/0504166}{arXiv:quant-ph/0504166}.

\bibitem{QIQC}
QIQCOP Zoo, \emph{Tough error models}, problem record
\href{https://qiqc-op.com/problem/op_b1b41f737f9b4aeb/}{op\_b1b41f737f9b4aeb},
version dated 4 September 2026; accessed 13 November 2026.

\bibitem{KL97}
E.~Knill and R.~Laflamme, Theory of quantum error-correcting codes,
\emph{Physical Review A} \textbf{55}, 900--911 (1997),
\href{https://doi.org/10.1103/PhysRevA.55.900}{doi:10.1103/PhysRevA.55.900}.

\bibitem{KLV00}
E.~Knill, R.~Laflamme, and L.~Viola,
Theory of quantum error correction for general noise,
\emph{Physical Review Letters} \textbf{84}, 2525--2528 (2000),
\href{https://doi.org/10.1103/PhysRevLett.84.2525}{doi:10.1103/PhysRevLett.84.2525};
\href{https://arxiv.org/abs/quant-ph/9908066}{arXiv:quant-ph/9908066}.

\bibitem{Reimpell}
M.~Reimpell, \emph{Quantum Information and Convex Optimization},
doctoral dissertation, Technische Universit\"at Braunschweig (2008),
Section~3.3.4,
\href{https://www.osti.gov/etdeweb/servlets/purl/21161752}{archived dissertation}.

\bibitem{Weaver}
N.~Weaver, The ``quantum'' Tur\'an problem for operator systems,
\emph{Pacific Journal of Mathematics} \textbf{301}, 335--349 (2019),
\href{https://doi.org/10.2140/pjm.2019.301.335}{doi:10.2140/pjm.2019.301.335};
\href{https://arxiv.org/abs/1802.07394}{arXiv:1802.07394}.

\bibitem{AK}
A.~Allen and A.~Kornell, The quantum Ramsey numbers $QR(2,k)$,
\emph{Linear and Multilinear Algebra} \textbf{73}, 3959--3964 (2025),
\href{https://doi.org/10.1080/03081087.2025.2578529}{doi:10.1080/03081087.2025.2578529};
\href{https://arxiv.org/abs/2507.04128v2}{arXiv:2507.04128v2}.

\bibitem{LiTam}
C.-K.~Li and B.-S.~Tam, A note on extreme correlation matrices,
\emph{SIAM Journal on Matrix Analysis and Applications} \textbf{15},
903--908 (1994),
\href{https://doi.org/10.1137/S0895479892240683}{doi:10.1137/S0895479892240683}.

\bibitem{Gupta}
A.~Gupta, B.~Sun, X.~He, and B.~Zeng,
\emph{Variance Geometry of Exact Pauli-Detecting Codes:
Continuous Landscapes Beyond Stabilizers} (2026),
\href{https://arxiv.org/abs/2604.21800v1}{arXiv:2604.21800v1}.

\bibitem{Sun}
B.~Sun, N.~Cao, Y.~T.~Poon, and B.~Zeng,
\emph{Geometry of Knill--Laflamme Coefficients for Pauli Error Detection}
(2026), \href{https://arxiv.org/abs/2610.01992v1}{arXiv:2610.01992v1}.

\bibitem{CKZ}
M.-D.~Choi, D.~W.~Kribs, and K.~\.{Z}yczkowski,
Higher-rank numerical ranges and compression problems,
\emph{Linear Algebra and its Applications} \textbf{418}, 828--839 (2006),
\href{https://doi.org/10.1016/j.laa.2006.03.019}{doi:10.1016/j.laa.2006.03.019};
\href{https://arxiv.org/abs/math/0511278}{arXiv:math/0511278}.
\end{thebibliography}
\end{document}
